\newpage
\section{Componente Frontend: Modulo Export}

\subsection{Descrizione}
Il modulo Export gestisce l'\textit{application logic} dedicata all'estrazione, salvataggio e stampa dei documenti contenuti nei pacchetti di archiviazione. Il modulo coordina le operazioni di trasferimento dei file dal sistema protetto (DIP) verso il \textit{file system} locale dell'utente o verso le periferiche di output. L'architettura segue il pattern Smart-Dumb Components per la gestione dei progressi di scaricamento e il pattern Facade per l'orchestrazione delle richieste IPC.

\subsubsection{Lista delle Classi e Strutture Dati}

\paragraph{ExportState}
Oggetto responsabile della gestione dello stato reattivo delle operazioni di output tramite Angular Signals. Traccia il ciclo di vita delle richieste di esportazione e la persistenza temporanea dei risultati.

\paragraph{Descrizione degli attributi della struttura:}
\begin{itemize}
\item \textbf{phase Signal}: segnale che rappresenta lo stato globale dell'operazione (IDLE, EXPORTING, COMPLETED, ERROR);
\item \textbf{progress Signal}: valore numerico (0-100) che indica la percentuale di avanzamento dell'operazione di esportazione corrente;
\item \textbf{lastResult Signal<ExportResult | null>}: memorizza l'esito dell'ultima operazione (es. percorso di salvataggio, numero di file estratti);
\item \textbf{error Signal<AppError | null>}: contiene l'eventuale errore riscontrato durante la comunicazione IPC o l'accesso al disco;
\item \textbf{loading Signal}: indica se il modulo è in attesa di una risposta dal processo Main per un'operazione di I/O.
\end{itemize}

\paragraph{Descrizione dei metodi invocabili dalla struttura:}
\begin{itemize}
\item \textbf{setExporting()}: imposta la fase su EXPORTING e resetta il progresso;
\item \textbf{updateProgress(value: number)}: aggiorna il segnale del progresso durante i download multipli;
\item \textbf{setSuccess(result: ExportResult)}: imposta lo stato su COMPLETED e memorizza i dettagli dell'esito;
\item \textbf{setError(error: AppError)}: registra un fallimento critico (es. permessi disco negati);
\item \textbf{reset()}: riporta il modulo allo stato IDLE.
\end{itemize}

\paragraph{ExportFacade}
Implementa la \textit{business logic} di coordinamento per le funzionalità di output. Funge da interfaccia unificata per i componenti della UI (\textit{ExportPage}, \textit{ExportProgress}) delegando le chiamate di sistema all'infrastruttura. Implementa l'interfaccia \textit{IExportFacade}.

\paragraph{Descrizione degli attributi della struttura:}
\begin{itemize}
\item \textbf{exportState ExportState}: istanza del servizio di gestione dello stato reattivo delle esportazioni;
\item \textbf{ipcGateway ExportIpcGateway}: istanza del gateway incaricato di dialogare con i canali di download di Electron.
\end{itemize}

\paragraph{Descrizione dei metodi invocabili dalla struttura:}
\begin{itemize}
\item \textbf{exportDocument(documentId string) Promise}: avvia la procedura di salvataggio per un singolo documento. Gestisce l'apertura della finestra di dialogo nativa tramite gateway e monitora l'esito;
\item \textbf{exportMultiple(documentIds string[]) Promise}: coordina l'estrazione massiva di documenti, aggiornando progressivamente lo stato di avanzamento per ogni file processato;
\item \textbf{printDocument(documentId string) Promise}: invoca il servizio di sistema per la stampa diretta del documento selezionato;
\item \textbf{cancelOperation() void}: tenta di interrompere un'operazione di esportazione in corso tramite segnale IPC.
\end{itemize}

\paragraph{Descrizione dei segnali esposti:}
\begin{itemize}
\item \textbf{progress}: espone l'avanzamento percentuale per la UI;
\item \textbf{phase}: permette alla UI di reagire al cambio di stato (es. mostrando lo spinner o il messaggio di successo);
\item \textbf{lastResult}: fornisce i metadati dell'operazione conclusa.
\end{itemize}

\subsection{Interfacce}

\subsubsection{IExportFacade}
Definisce il contratto per la Business Logic del modulo Export. Stabilisce i metodi necessari per le azioni di output e i segnali per il monitoraggio dei processi asincroni.

\paragraph{Descrizione dei metodi dell'interfaccia:}
\begin{itemize}
\item \textbf{exportDocument(id: string) Promise}: contratto per l'estrazione singola;
\item \textbf{exportMultiple(ids: string[]) Promise}: contratto per l'estrazione di gruppo;
\item \textbf{printDocument(id: string) Promise}: contratto per l'invio alla periferica di stampa.
\end{itemize}

\subsubsection{IExportChannel}
Interfaccia di contratto per il layer di infrastruttura. Definisce le firme per le chiamate IPC dedicate ai flussi di file tra il sandbox di Electron e il sistema operativo.

\paragraph{Descrizione dei metodi dell'interfaccia:}
\begin{itemize}
\item \textbf{saveFile(id: string) Promise}: definisce la richiesta di apertura del \textit{Save Dialog} nativo;
\item \textbf{print(id: string) Promise}: definisce la richiesta di stampa;
\item \textbf{openFolder(path: string) void}: comando per aprire la cartella di destinazione nel file manager di sistema.
\end{itemize}

\subsubsection{ExportIpcGateway}
Adapter di infrastruttura che implementa la comunicazione asincrona verso i canali \texttt{file:*} di Electron. Utilizza il bridge fornito dal contesto di esecuzione desktop.

\paragraph{Descrizione dei metodi invocabili dalla struttura:}
\begin{itemize}
\item \textbf{saveFile(id string)}: invoca il canale \texttt{file:download-single} passando l'identificativo del documento e riceve il percorso finale scelto dall'utente;
\item \textbf{print(id string)}: esegue l'invocazione del canale \texttt{file:print-document};
\item \textbf{requestMultiExport(ids string[])}: inizializza una sessione di esportazione massiva sul canale \texttt{file:download-batch}.
\end{itemize}

\subsection{Dipendenze}
\begin{itemize}
\item \textbf{Angular Signals}: gestione della reattività del progresso e degli stati di errore;
\item \textbf{Electron IPC Bridge}: per l'accesso alle \textit{Native APIs} del sistema operativo (File System e Print Manager);
\item \textbf{Domain Models/DTOs}: definizioni condivise per rappresentare i risultati dell'esportazione (\textit{ExportResult}) e i metadati di trasporto.
\end{itemize}

\subsection{Tracciamento dei requisiti}
La tabella seguente mappa i requisiti funzionali obbligatori relativi all'Export sulle componenti realizzate.

\begin{table}[H]
    \centering
    \begin{tabular}{|p{2.5cm}|p{6.5cm}|p{5.5cm}|}
        \hline
        \rowcolor{zpusgreen!30}
        \textbf{Codice} & \textbf{Descrizione Sintetica} & \textbf{Componente Software} \\
        \hline
        \textbf{R-178-F-Ob} & Salvataggio di un documento in locale. & \textit{ExportFacade, ExportIpcGateway} \\
        \hline
        \textbf{R-179-F-Ob} & Salvataggio multiplo di documenti. & \textit{ExportFacade, ExportProgress} \\
        \hline
        \textbf{R-181-F-Ob} & Stampa di un documento. & \textit{ExportFacade, ExportIpcGateway} \\
        \hline
        \textbf{R-224-F-Ob} & Selezione della cartella di destinazione. & \textit{IExportChannel, ExportPage} \\
        \hline
        \textbf{R-226-F-Ob} & Protezione della cartella sorgente del DIP. & \textit{ExportFacade (Business Logic)} \\
        \hline
    \end{tabular}
    \caption{Tracciamento Requisiti Funzionali - Modulo Export}
\end{table}